% ===================================================================
%  CUADRO N.º 4 — Edad Madura y Vejez.
%  Traduction 2026 d'apres texto/io, controlee sur texto/fr et
%  texto/en. Consigne : voir texto/es/10-cuadro-01.tex.
% ===================================================================

%%K t04-apar-1 apar
\VUsaut{4.0mm}
\VUtitre{15.0pt}{60}{CUADRO N.º 4}
\VUsaut{1.6mm}
\VUtitre{11.4pt}{0}{Edad Madura y Vejez.}
\VUsaut{3.2mm}

%%K t04-tit-1 sub
\VUcentre{11.4pt}{0}{\textit{Primera escena.}}
\VUsaut{1.6mm}
\VUtitre{11.4pt}{0}{El Banquete de Boda.}
\VUsaut{2.6mm}

%%K t04-tit-2 sub
\VUpk{10.2pt}{40}{El otoño.}
\VUsaut{3.0mm}

%%K t04-c1-01-1 p
1. --- Los jóvenes han crecido. Uno de ellos, Juan, se casa. Asistimos al
\VUgras{banquete de boda}, que reúne alrededor de la misma mesa a las
familias de los dos esposos.

%%K t04-c1-02-1 p
2. --- Estamos en \VUgras{otoño}, en el mes de \VUgras{septiembre}. El viento
frío de \VUgras{octubre} y de \VUgras{noviembre} no ha despojado todavía al
campo de su verde ropaje. El aire de la tarde es tibio y perfumado; por
eso la cena se sirve bajo la \VUgras{galería} \textsuperscript{(8)} que da al
jardín.

%%K t04-c1-03-1 p
3. --- A lo lejos, en la \VUgras{colina} \textsuperscript{(3)}, está la casa de los
padres de Juan, una elegante \VUgras{casa de campo} \textsuperscript{(1)} escondida
entre el verdor; hermosos viñedos, que dan un vino excelente, cubren la
ladera del cerro. El viñador los cultiva y los cuida.

%%K t04-c1-04-1 p
4. --- Por encima de las viñas pasa un bando de \VUgras{perdices}
\textsuperscript{(2)}, espantadas por la llegada del \VUgras{guarda de caza}
\textsuperscript{(5)}. Este, con la \VUgras{escopeta} bajo el brazo y el
\VUgras{morral} en bandolera, bate los \VUgras{matorrales} \textsuperscript{(4)} con
su \VUgras{perro} \textsuperscript{(6)}. Vigila la finca e impide que los furtivos
acaben con la caza.

%%K t04-c1-05-1 p
5. --- Fuera de la \VUgras{galería} trepa una \VUgras{parra}, de la que cuelgan
pesados \VUgras{racimos de uva} \textsuperscript{(7)}.

%%K t04-c1-06-1 p
6. --- Las bonitas flores de otoño que adornan la \VUgras{mesa}
\textsuperscript{(16)} son \VUgras{crisantemos} \textsuperscript{(9)}; el \VUgras{padre}
\textsuperscript{(10)} de la novia se ha puesto uno en el ojal del frac.

%%K t04-c1-07-1 p
7. --- La comida toca a su fin. El \VUgras{criado} \textsuperscript{(11)} empieza a
traer los postres y va a retirar las distintas piezas del cubierto ---
\VUgras{cucharas} \textsuperscript{(14)}, \VUgras{tenedores} \textsuperscript{(12)},
\VUgras{cuchillos} \textsuperscript{(13)}, \VUgras{fuentes} \textsuperscript{(15)} --- que ya
no hacen falta.

%%K t04-tit-3 sub
\VUpk{10.2pt}{40}{La Familia.}
\VUsaut{3.0mm}

%%K t04-c1-08-1 p
8. --- Todos parecen alegres, y la sonrisa con que la \VUgras{madre}
\textsuperscript{(17)} del novio mira a sus hijos prueba su felicidad.

%%K t04-c1-09-1 p
9. --- Esta boda une a dos familias ricas de la comarca. Juan, el
\VUgras{marido} \textsuperscript{(18)}, es \VUgras{hijo} de un gran propietario y se
convierte en el \VUgras{yerno} de un gran industrial.

%%K t04-c1-10-1 p
10. --- Los \VUgras{esposos} son jóvenes y sin embargo llevaban mucho tiempo
\VUgras{prometidos}. La \VUgras{joven esposa} \textsuperscript{(19)}, Marta, está
encantadora con su \VUgras{vestido blanco} adornado con flores de azahar.

%%K t04-c1-11-1 p
11. --- Su \VUgras{suegro} \textsuperscript{(20)} está muy contento de tener una
\VUgras{nuera} tan amable, y la \VUgras{suegra} \textsuperscript{(21)} de Juan sonríe
al ver a su \VUgras{hija} tan hermosa.

%%K t04-c1-12-1 p
12. --- En el extremo de la mesa se sientan los demás \VUgras{invitados}
\textsuperscript{(22)}: \VUgras{parientes, amigos, primos y primas}. Un
\VUgras{maître} \textsuperscript{(23)}, con la servilleta bajo el brazo, los sirve
con diligencia.

%%K t04-tit-4 sub
\VUpk{10.2pt}{40}{Los Amigos.}
\VUsaut{3.0mm}

%%K t04-c1-13-1 p
13. --- Al lado derecho de la mesa, he aquí a una \VUgras{señorita}
\textsuperscript{(24)}, \VUgras{amiga} de la novia, y luego al \VUgras{hermano} de
esta, que es su \VUgras{padrino de boda} \textsuperscript{(25)}. Charla con su
\VUgras{dama de honor} \textsuperscript{(26)}, hermana del novio, que por esta boda
se convierte en \VUgras{cuñada} de Marta. Es una joven encantadora. Lleva
hermosas \VUgras{joyas}, un \VUgras{collar de perlas} y una \VUgras{pulsera} de
oro.

%%K t04-c1-14-1 p
14. --- Cerca de ellos, el señor que está de pie y levanta su copa es un
\VUgras{amigo} \textsuperscript{(27)} de la familia. Es uno de los
\VUgras{testigos del novio}. Ofrece un \VUgras{brindis}, bebe a la
\VUgras{salud} de los recién casados y les desea mucha felicidad. Va de
etiqueta, con frac y \VUgras{zapatos de charol} \textsuperscript{(28)}. Acompaña a
la \VUgras{abuela} \textsuperscript{(29)}, que es \VUgras{viuda}.

%%K t04-c1-15-1 p
15. --- El joven de \VUgras{frac} \textsuperscript{(31)}, del fino bigote negro, es
el \VUgras{cuñado} \textsuperscript{(30)} de Juan. Es un ingeniero distinguido. Está
al lado de una \VUgras{joven señora} \textsuperscript{(33)} muy elegante, la
\VUgras{hermana mayor} de Marta y madre de la preciosa niña que viene, del
brazo de su primo, a ofrecerle una flor y a \VUgras{darle} las
\VUgras{buenas noches}. Esta señora lleva unos \VUgras{pendientes} muy
hermosos y un elegante \VUgras{peinado}.

%%K t04-c1-16-1 p
16. --- El primito, tan gracioso con su \VUgras{blusón} \textsuperscript{(36)} y su
\VUgras{pantalón} \textsuperscript{(37)} abombado, da mucha lástima. Es
\VUgras{huérfano} \textsuperscript{(35)}. Perdió muy joven a su padre y a su madre.
Su tío es su \VUgras{tutor} \textsuperscript{(38)} y se quedó soltero para criar al
pobre niño. Es un hombre excelente. Está algo calvo y es muy miope; lleva
\VUgras{quevedos}.

%%K t04-tit-5 sub
\VUsaut{2.4mm}
\VUpk{10.2pt}{40}{El Postre.}
\VUsaut{3.0mm}

%%K t04-c1-17-1 p
17. --- Detrás de los comensales, en una mesa cubierta con un
\VUgras{mantel}, está preparado el \VUgras{postre} \textsuperscript{(39)}. En una
copa grande, he aquí fruta: \VUgras{melocotones} \textsuperscript{(40)},
\VUgras{peras} \textsuperscript{(41)}, \VUgras{manzanas} \textsuperscript{(42)},
\VUgras{albaricoques} \textsuperscript{(43)} y \VUgras{uvas} \textsuperscript{(44)}. Al lado,
\VUgras{piñas} \textsuperscript{(45)}, un hermoso \VUgras{pastel} \textsuperscript{(46)},
\VUgras{botellas de licor} \textsuperscript{(48)} y \VUgras{champán} \textsuperscript{(49)}, y
también pilas de \VUgras{platos} \textsuperscript{(47)} limpios.

%%K t04-c1-18-1 p
18. --- El \VUgras{sumiller} \textsuperscript{(50)} descorcha una \VUgras{botella}
\textsuperscript{(52)} de vino añejo con un \VUgras{sacacorchos} \textsuperscript{(51)}. El
\VUgras{corcho} \textsuperscript{(53)} parece muy difícil de sacar. Este sumiller
lleva una \VUgras{servilleta} \textsuperscript{(54)} bajo el brazo.

%%K t04-c1-19-1 p
19. --- Después de la cena, los señores pasarán al salón de fumar y las
señoras irán a arreglarse para el baile.

%%K t04-tit-6 sub
\VUsaut{5.0mm}
\VUcentre{11.4pt}{0}{\textit{Segunda escena.}}
\VUsaut{1.6mm}
\VUpk{11.4pt}{40}{El Cumpleaños del abuelo.}
\VUsaut{2.6mm}

%%K t04-tit-7 sub
\VUpk{10.2pt}{40}{La Visita.}
\VUsaut{3.0mm}

%%K t04-c2-01-1 p
1. --- Han pasado diez años. Los padres de Juan se han hecho viejos y
quieren mucho a sus tres nietos, dos \VUgras{nietos} \textsuperscript{(2)} y una
\VUgras{nieta} \textsuperscript{(3)}. Como hoy es el \VUgras{cumpleaños del abuelo},
los hijos vienen a felicitarlo.

%%K t04-c2-02-1 p
2. --- El pequeño Pedro es todavía muy niño; solo tiene tres años. Ve que
se acerca el \VUgras{gato} \textsuperscript{(1)}. Quiere acariciar su hermoso
\VUgras{pelo} sedoso y su larga \VUgras{cola}; no piensa que, bajo esas
graciosas \VUgras{patas}, se esconden feas uñas puntiagudas que bien
podrían arañarlo.

%%K t04-c2-03-1 p
3. --- Su hermana Gabriela, que le da la mano, es mucho más seria, y
Marcelo, con su gran \VUgras{abrigo} \textsuperscript{(4)} y su \VUgras{cinturón}
\textsuperscript{(5)} de cuero, parece ya un hombrecito, a pesar del pantalón
corto que deja ver sus \VUgras{medias} \textsuperscript{(6)}.

%%K t04-c2-04-1 p
4. --- Su padre se ha puesto su grueso \VUgras{gabán} \textsuperscript{(7)} de
invierno con \VUgras{cuello de piel} \textsuperscript{(8)} y, en el \VUgras{bolsillo}
\textsuperscript{(9)}, lleva un pañuelo. Su mujer lleva su sombrero de fieltro
adornado con \VUgras{plumas} \textsuperscript{(10)}, su \VUgras{chaqueta}
\textsuperscript{(11)}, su \VUgras{boa de piel} \textsuperscript{(12)} y su \VUgras{manguito}
\textsuperscript{(13)}.

%%K t04-c2-05-1 p
5. --- Hace mucho frío, porque reina el invierno. La \VUgras{nieve}
\textsuperscript{(19)} ha caído todo el día y los tejados de las casas están
blancos. Con la \VUgras{noche}, el frío se hace aún más vivo. En el cielo,
la \VUgras{luna} \textsuperscript{(17)} y las \VUgras{estrellas} \textsuperscript{(18)}
brillan y atraviesan la \VUgras{oscuridad}.

%%K t04-tit-8 sub
\VUsaut{4.2mm}
\VUpk{10.2pt}{40}{La Enfermedad.}
\VUsaut{3.0mm}

%%K t04-c2-06-1 p
6. --- El pobre \VUgras{ciego} \textsuperscript{(20)} que cruza la plaza delante de
la \VUgras{farmacia} \textsuperscript{(16)}, guiado por su \VUgras{caniche}
\textsuperscript{(21)}, es muy desdichado. Justina, la \VUgras{criada}
\textsuperscript{(14)}, trae de la cocina un \VUgras{tazón} \textsuperscript{(15)} de
\VUgras{infusión} muy caliente y bien azucarada.

%%K t04-c2-07-1 p
7. --- La \VUgras{abuela} \textsuperscript{(23)}, que quiere ir a buscar sus
\VUgras{gafas} \textsuperscript{(26)} a la mesa para leer la \VUgras{receta}
\textsuperscript{(27)} que el \VUgras{médico} \textsuperscript{(25)} acaba de escribir, se
levanta de su sillón apoyándose en su \VUgras{bastón} \textsuperscript{(24)}.

%%K t04-c2-08-1 p
8. --- A menudo sufre pequeñas \VUgras{indisposiciones}, \VUgras{mareos},
\VUgras{resfriados} y \VUgras{dolores de muelas}; tiene las piernas débiles y
la vista ya no es muy buena. A pesar de ello, se burla de buena gana de
su viejo \VUgras{médico}, que siempre quiere recetarle una \VUgras{poción}
\textsuperscript{(28)}. Hoy está muy contenta de tener a su joven familia a su
alrededor.

%%K t04-c2-09-1 p
9. --- Se está a gusto en la habitación bien cerrada. En la
\VUgras{chimenea} \textsuperscript{(29)} de mármol arde un \VUgras{fuego espléndido}
\textsuperscript{(30)}, y uno se siente feliz con la dulce \VUgras{luz} de la
\VUgras{lámpara} \textsuperscript{(31)} que acaban de encender.

%%K t04-tit-9 sub
\VUsaut{4.2mm}
\VUpk{10.2pt}{40}{El Abuelo.}
\VUsaut{3.0mm}

%%K t04-c2-10-1 p
10. --- Hasta el \VUgras{abuelo} \textsuperscript{(33)} se olvida de que ha estado
enfermo, de que su \VUgras{reuma} lo tiene clavado en su \VUgras{sillón}
\textsuperscript{(34)} y de que, ayer mismo, tenía \VUgras{fiebre y desmayos}. Ya
no se acuerda de que el invierno pasado tuvo una \VUgras{pulmonía} y de que
una \VUgras{tos} ronca y penosa lo hacía sufrir mucho.

%%K t04-c2-10-2 p
Y sin embargo se restableció con mucha dificultad. Ahora está algo mejor.
Pero la pierna, que apoya sobre un \VUgras{cojín} \textsuperscript{(38)} puesto en
un pequeño \VUgras{taburete} \textsuperscript{(39)}, todavía le duele.

%%K t04-c2-11-1 p
11. --- Cuando está bien envuelto en su cálida \VUgras{bata} \textsuperscript{(35)}
de gruesos \VUgras{botones} \textsuperscript{(36)} de nácar, y cuando tiene a su
lado, para leerle el periódico, a la pequeña \VUgras{huérfana}
\textsuperscript{(40)} con su \VUgras{cofia} blanca, su \VUgras{vestido}
\textsuperscript{(42)} y su \VUgras{esclavina} \textsuperscript{(41)} azules, no se queja.

%%K t04-c2-12-1 p
12. --- Todos los años, cuando el \VUgras{calendario} \textsuperscript{(43)} trae de
nuevo la \VUgras{fecha} del primero de \VUgras{diciembre}, espera con
impaciencia a sus \VUgras{nietos}, porque sabe que no olvidarán esa
\VUgras{fecha} importante.

%%K t04-c2-13-1 p
13. --- Recuerda con emoción los tiempos lejanos en que él mismo iba a
felicitar al glorioso \VUgras{mariscal} del Imperio, cuyo \VUgras{retrato}
\textsuperscript{(44)} cuelga de la pared, y que es el \VUgras{bisabuelo} de los
niños.
\VUsaut{6.0mm}
\VUfilet{22mm}
